\subsection{消元定理}

命题 9. 设 \(\mathbf{S}_1\) 与 \(\mathbf{S}_2\) 是两个不相交的集合， \(d\) 是 \(\mathbf{S}_1 \times \mathbf{S}_2\) 到 \(\mathrm{L}(\mathbf{S}_2)\) 中的一个映射。设 \(\mathfrak{g}\) 是 \(\mathrm{L}(\mathbf{S}_1 \cup \mathbf{S}_2)\) 对由元素 \([s_1, s_2] - d(s_1, s_2)\) （满足 \(s_1 \in \mathbf{S}_1\) 与 \(s_2 \in \mathbf{S}_2\) ）生成的理想所作的商李代数；设 \(\psi\) 是 \(\mathrm{L}(\mathbf{S}_1 \cup \mathbf{S}_2)\) 到 \(\mathfrak{g}\) 上的典范映射。

\begin{enumerate}
\def\labelenumi{(\alph{enumi})}
\item
  对 \(i = 1,2\) ， \(\phi_{i}\) （即 \(\psi\) 在 \(S_{i}\) 上的限制）可以扩张为 \(\mathbf{L}(S_i)\) 到一个子代数 \(\mathfrak{a}_i\) （包含于 \(\mathfrak{g}\) ）上的同构。
\item
  \(\mathfrak{g} = \mathfrak{a}_1 + \mathfrak{a}_2, \mathfrak{a}_1 \cap \mathfrak{a}_2 = \{0\}\) ，并且 \(\mathfrak{a}_2\) 是 \(\mathfrak{g}\) 的一个理想。
\end{enumerate}

对 \(i = 1,2\) ，用 \(\psi_{i}\) 表示 \(\mathbf{L}(\mathbf{S}_i)\) 到 \(\mathfrak{g}\) 中的扩张 \(\phi_i\) 的同态，并用 \(\mathfrak{a}_i\) 表示它的像。显然 \(\phi_i(\mathbf{S}_i)\) 生成 \(\mathfrak{a}_i\) 。

设 \(s_1 \in S_1\) ；我们记 \(D = \operatorname{ad} \phi_1(s_1)\) 。 \(D\) 是 \(g\) 的一个导子，它把 \(\phi_2(S_2)\) 映到 \(a_2\) 之中，依照关系式

\[
[ \phi_ {1} (s _ {1}), \phi_ {2} (s _ {2}) ] = \psi_ {2} (d (s _ {1}, s _ {2})) \quad \text { for } s _ {2} \in S _ {2};
\]

由于 \(\mathfrak{a}_2\) 作为 \(\mathfrak{g}\) 的子代数由 \(\phi_2(S_2)\) 生成，因此 \(D(\mathfrak{a}_2) \subset \mathfrak{a}_2\) 。由那些 \(x \in \mathfrak{g}\) （它们使 ad \(x\) 保持 \(\mathfrak{a}_2\) 不变）组成的集合是 \(\mathfrak{g}\) 的一个李子代数，由上所述它包含 \(\phi_1(S_1)\) ，从而也包含 \(\mathfrak{a}_1\) 。因此

\[
\left[ \mathfrak {a} _ {1}, \mathfrak {a} _ {2} \right] \subset \mathfrak {a} _ {2}.\tag{23}
\]

因此 \(a_1 + a_2\) 是 \(g\) 的一个李子代数，并且由于它包含生成集 \(\phi_1(S_1) \cup \phi_2(S_2)\) ，

\[
a _ {1} + a _ {2} = g.\tag{24}
\]

对一切 \(s_1 \in S_1\) ，存在导子 \(D_{s_1}\) （\(L(S_2)\) 上的），使得

\[
\mathrm{D} _ {s _ {1}} (s _ {2}) = d (s _ {1}, s _ {2})
\]

对一切 \(s_2\) （属于 \(S_2\) ）成立 (no. 8， 命题 8 的推论)。映射 \(s_1 \mapsto D_{s_1}\) 可以扩张为一个同态 \(D\) （从 \(L(S_1)\) 到 \(L(S_2)\) 的导子李代数中）。设 \(\mathfrak{h}\) 是 \(L(S_1)\) 与 \(L(S_2)\) 的对应于 \(D\) 的半直积 (第 I 章， § 1， no. 8)。作为模， \(\mathfrak{h}\) 等于 \(L(S_1) \times L(S_2)\) ，特别地

\[
[ (s _ {1}, 0), (0, s _ {2}) ] = (0, d (s _ {1}, s _ {2}))\tag{25}
\]

其中 \(s_{1} \in S_{1}\) 与 \(s_{2} \in S_{2}\) 。

由 (25) 我们推出存在一个同态 \(f\) （从 \(\mathfrak{g}\) 到 \(\mathfrak{h}\) 中），使得 \(f(\phi_1(s_1)) = (s_1, 0)\) ，并且 \(f(\phi_2(s_2)) = (0, s_2)\) （对 \(s_1 \in S_1\) 与 \(s_2 \in S_2\) ）。我们立即推出关系式

\[
f (\psi_ {1} (a _ {1}) + \psi_ {2} (a _ {2})) = (a _ {1}, a _ {2})\tag{26}
\]

其中 \(a_1 \in \mathrm{L}(\mathrm{S}_1)\) 与 \(a_2 \in \mathrm{L}(\mathrm{S}_2)\) 。

关系式 (26) 表明 \(\psi_{1}\) 与 \(\psi_{2}\) 都是单射，并且 \(a_{1} \cap a_{2} = \{0\}\) 。于是公式 (23) 与 (24) 蕴涵本命题。

命题 10 (消元定理). 设 X 是一个集合，S 是 X 的一个子集，T 是由这样的序列 \((s_1, \ldots, s_n, x)\) 组成的集合：\(n \geqslant 0\) ， \(s_1, \ldots, s_n\) 属于 S 且 \(x\) 属于 \(\mathbf{X} - \mathbf{S}\) 。

\begin{enumerate}
\def\labelenumi{(\alph{enumi})}
\item
  模 \(\mathbf{L}(\mathbf{X})\) 是子代数 \(\mathbf{L}(\mathbf{S})\) （\(\mathbf{L}(\mathbf{X})\) 中的）与理想 \(\mathfrak{a}\) （\(\mathbf{L}(\mathbf{X})\) 中由 \(\mathbf{X} - \mathbf{S}\) 生成的）的直和。
\item
  存在一个李代数同构 \(\phi\) （从 L(T) 到 \(\mathfrak{a}\) 上），它把 \((s_1, \ldots, s_n, x)\) 映为 (ad \(s_1 \circ \cdots \circ\) ad \(s_n\) ) ( \(x\) )。
\end{enumerate}

设 \(g\) 是按命题 9 的方式构造的李代数，给定

\[
\mathrm{S} _ {1} = \mathrm{S}, \quad \mathrm{S} _ {2} = \mathrm{T}, \quad d (s, t) = (s, s _ {1}, \dots , s _ {n}, x) \in \mathrm{T} \subset \mathrm{L} (\mathrm{T})
\]

其中 \(t = (s_1, \ldots, s_n, x)\) 属于 \(T\) 且 \(s \in S_1\) 。我们把 \(L(S)\) 与 \(L(T)\) 和它们在 \(g\) 中的典范像等同 (命题 9 (a))。

设 \(\psi\) 是映射 \((s_1, \ldots, s_n, x) \mapsto (\operatorname{ad} s_1 \circ \cdots \circ \operatorname{ad} s_n)(x)\) （从 T 到 \(\mathrm{L}(\mathbf{X})\) 中）。显然 \(\psi(d(s, t)) = [s, \psi(t)]\) （对 \(s \in S\) 与 \(t \in T\) ），因而存在一个同态 \(\alpha: g \to L(X)\) ，它在 \(S\) 上的限制是恒同映射，而在 \(T\) 上的限制是 \(\psi\) 。而 \(X - S \subset T\) ，从而存在一个同态 \(\beta: L(X) \to g\) ，它在 \(X = S \cup (X - S)\) 上的限制是恒同映射。

我们证明 \(\alpha\) 是一个同构，而 \(\beta\) 是其逆同构。由于 \(\psi(x) = x\) 对每个 \(x\) （属于 \(\mathbf{X} - \mathbf{S}\) ）成立，我们看出 \(\alpha \circ \beta\) 在 \(\mathbf{X}\) 上与恒同映射一致，从而 \(\alpha \circ \beta = \mathrm{Id}_{\mathrm{L}(\mathbf{X})}\) 。另一方面，按构造， \([s, t] = d(s, t)\) 属于 \(\mathfrak{g}\) （对 \(s \in \mathbb{S}\) 与 \(t \in \mathbb{T}\) ）；由此推出 \(t = (s_1, \ldots, s_n, x)\) 在 \(\mathfrak{g}\) 中等于 \((\operatorname{ad} s_1 \circ \cdots \circ \operatorname{ad} s_n)(x)\) ，从而 \(t = \beta(\alpha(t))\) 。由于 \(\beta(\alpha(s)) = s\) （对 \(s \in \mathbb{S}\) ），并且 \(\mathbb{S} \cup \mathbb{T}\) 生成 \(\mathfrak{g}, \beta \circ \alpha = \mathrm{Id}_{\mathfrak{g}}\) 。

由于 \(\alpha\) 是 \(g\) 到 \(L(X)\) 上的一个同构，命题 9 表明 \(\alpha\) 在 \(\mathrm{L(T)}\) 上的限制是一个同构 \(\phi\) ，它把 \(\mathrm{L(T)}\) 映到一个理想 \(\mathfrak{b}\) （\(\mathrm{L(X)}\) 中的）上，使得模 \(\mathrm{L(X)}\) 是 \(\mathrm{L(S)}\) 与 \(\mathfrak{b}\) 的直和。显然

\[
\phi (s _ {1}, \dots , s _ {n}, x) = (\operatorname{ad} s _ {1} \circ \dots \circ \operatorname{ad} s _ {n}) (x)
\]

其中 T 中的 \((s_1, \ldots, s_n, x)\) 。

因此 \(\phi(T) \subset \mathfrak{a}\) ，而 \(\mathfrak{b} \subset \mathfrak{a}\) ，因为 \(\phi(T)\) 生成子代数 \(\mathfrak{b}\) （\(L(X)\) 中的）。但 \(\mathfrak{b}\) 是一个理想且 \(X - S \subset \phi(T) \subset \mathfrak{b}\) ，由此得 \(\mathfrak{a} \subset \mathfrak{b}\) 。

推论. 设 \(y \in \mathbf{X}\) 。自由李代数 \(\mathbf{L}(\mathbf{X})\) 是自由子模 \(\mathbf{K}.y\) 与以 \(((\operatorname{ad}y)^n.z)\) 组成的族（其中 \(n \geqslant 0\) 且 \(z \in \mathbf{X} - \{y\}\) ）为基本族的李子代数的直和。

只需在命题 10 中取 \(S = \{y\}\) 。

